\documentclass[12pt]{article}
\usepackage{ctex}
\usepackage{amsmath, amssymb}
\usepackage{geometry}
\geometry{a4paper, margin=2.5cm}
\usepackage{enumitem}
\title{第 3 章\ 功和能\ 例题答案}
\author{}
\date{}

\begin{document}
\maketitle

\begin{enumerate}[label=\textbf{例\arabic*.}, leftmargin=*]

% === 例 1 ===
\item 由 $\dfrac{F}{m} = 6t = \dfrac{\mathrm{d}v}{\mathrm{d}t}$，得 $v = 3t^2$。又 $v = \dfrac{\mathrm{d}x}{\mathrm{d}t}$，故 $\mathrm{d}x = 3t^2\mathrm{d}t$。
\[
  A = \int_0^x F\,\mathrm{d}x = \int_0^t 12t \cdot 3t^2\,\mathrm{d}t = \int_0^t 36t^3\,\mathrm{d}t = 144\,\mathrm{J}.
\]
\[
  P = \vec{F}\cdot\vec{v} = 12t \cdot 3t^2 = 288\,\mathrm{W} \quad (t=2\,\mathrm{s}).
\]

% === 例 2 ===
\item $v_x = 4t^2 \Rightarrow x = \dfrac{4}{3}t^3$；$v_y = 16 \Rightarrow y = 16t$。\\
$y=16 \Rightarrow t=1$；$y=32 \Rightarrow t=2$。\\
$F_x = m\dfrac{\mathrm{d}v_x}{\mathrm{d}t} = 80t$，$F_y = 0$。
\[
  A = \int_1^2 F_x\,\mathrm{d}x + \int F_y\,\mathrm{d}y = \int_1^2 80t \cdot 4t^2\,\mathrm{d}t = \int_1^2 320t^3\,\mathrm{d}t = 1200\,\mathrm{J}.
\]

% === 例 3 ===
\item 缓慢拉球时 $F - T\sin\theta = 0$，$T\cos\theta - mg = 0$，故 $F = mg\tan\theta$。\\
$\mathrm{d}s = L\,\mathrm{d}\theta$（$L$ 为绳长）：
\[
  A = \int_0^{\theta_0} F\cos\theta\,\mathrm{d}s = \int_0^{\theta_0} mgL\tan\theta\cos\theta\,\mathrm{d}\theta = mgL(1-\cos\theta_0).
\]

% === 例 4 ===
\item 设提升 $B$ 端使下端离地 $y$ 时被提起的部分长度为 $y/2$。提起部分质量 $\dfrac{M}{L}\cdot\dfrac{y}{2}$，其质心升高 $y/2$（在原下垂部分质心 $L/2$ 基础上再升 $y/2$）。\\
$\mathrm{d}A = -\dfrac{1}{2}\cdot\dfrac{M}{L}gy\,\mathrm{d}y$，积分：
\[
  A = \int_0^L -\frac{1}{2}\frac{M}{L}gy\,\mathrm{d}y = -\frac{1}{4}MgL.
\]

% === 例 5 ===
\item 离地面 $x$ 处，深 $\mathrm{d}x$ 的水层质量 $\mathrm{d}m = \rho S\,\mathrm{d}x$，将其提到地面（高度 $x$）的元功 $\mathrm{d}A = \rho S g x\,\mathrm{d}x$。
\[
  A = \int_H^{H+h} \rho S g x\,\mathrm{d}x = \frac{1}{2}\rho S g x^2\Big|_H^{H+h} = \rho g S H h + \rho g S \frac{h^2}{2}.
\]

% === 例 6 ===
\item $\mathrm{d}A = F\,\mathrm{d}s\cos\theta$。$\theta = Bs \Rightarrow \mathrm{d}s = \mathrm{d}\theta/B$。$\theta: 0 \to \pi/2$。
\[
  A = \int_0^{\pi/2} \frac{F\cos\theta}{B}\,\mathrm{d}\theta = \frac{F\sin\theta}{B}\Big|_0^{\pi/2} = \frac{F}{B}.
\]

% === 例 7 ===
\item $F_x = y$，$F_y = 3x^2$。\\
路径 $o \to b$：\\
$\mathrm{d}A = y\,\mathrm{d}x + 3x^2\,\mathrm{d}y$。\\
$o \to a$（$y=0$）：$A_1 = 0$。\\
$a \to b$（$x=2$, $\mathrm{d}x=0$）：$A_2 = \int_0^1 3\cdot 4\,\mathrm{d}y = 12$。故 $A_{oab} = 12\,\mathrm{J}$。\\
$o \to b$ 沿直线 $y = x/2$：$\mathrm{d}y = \mathrm{d}x/2$：
\[
  A = \int_0^2 \left(\frac{x}{2} + 3x^2\cdot\frac{1}{2}\right)\mathrm{d}x = \int_0^2 \left(\frac{x}{2} + \frac{3x^2}{2}\right)\mathrm{d}x = 1 + 4 = 5\,\mathrm{J}.
\]
说明 $\vec{F}$ 不是保守力。

% === 例 8 ===
\item $A_1 = mg\lambda_{\max}$（重力正功）；$A_2 = \displaystyle\int_0^{\lambda_{\max}}(-kx)\mathrm{d}x = -\frac{1}{2}k\lambda_{\max}^2$（弹性力负功）。\\
由动能定理（始末速度均为 0）：$mg\lambda_{\max} - \dfrac{1}{2}k\lambda_{\max}^2 = 0$，解得 $\lambda_{\max} = \dfrac{2mg}{k}$。

% === 例 9 ===
\item（1）设下垂长度 $y$ 时拉力 $\rho y g$，最大静摩擦 $\mu_0\rho(l - y)g$。临界 $y = b_0$：
\[
  \rho b_0 g - \mu_0\rho(l - b_0)g = 0 \Rightarrow b_0 = \frac{\mu_0}{1+\mu_0}l.
\]
当 $y > b_0$ 时链条开始滑动。\\
（2）以整个链条为对象：内力功之和为零。\\
重力做功 $A = \displaystyle\int_b^l \rho y g\,\mathrm{d}y = \frac{1}{2}\rho g(l^2 - b^2)$。\\
摩擦力做功 $A' = -\displaystyle\int_b^l \mu\rho(l - y)\,\mathrm{d}y = -\frac{1}{2}\mu\rho g(l - b)^2$。\\
由动能定理：$\dfrac{1}{2}\rho g(l^2 - b^2) - \dfrac{1}{2}\mu\rho g(l - b)^2 = \dfrac{1}{2}\rho l v^2$，解得
\[
  v = \sqrt{\frac{g}{l}(l^2 - b^2) - \frac{\mu g}{l}(l - b)^2}.
\]

% === 例 10 ===
\item（1）$\mathrm{d}A = F\,\mathrm{d}r\cos\theta = F R\,\mathrm{d}\varphi\sin\alpha$。由正弦定理 $\dfrac{\sin\alpha}{S} = \dfrac{\sin\varphi}{r}$，故 $\sin\alpha = \dfrac{S}{r}\sin\varphi$。\\
$\mathrm{d}A = m\mu r R\,\mathrm{d}\varphi\cdot\dfrac{S}{r}\sin\varphi = m\mu R S\sin\varphi\,\mathrm{d}\varphi$。
\[
  A = \int_0^\varphi m\mu R S\sin\varphi\,\mathrm{d}\varphi = -m\mu R S\cos\varphi\Big|_0^\varphi = m\mu R S(1 - \cos\varphi).
\]
（2）由动能定理：$m\mu R S(1 - \cos\varphi) = \dfrac{1}{2}m R^2\left(\dfrac{\mathrm{d}\varphi}{\mathrm{d}t}\right)^2$。\\
解出 $\dfrac{\mathrm{d}\varphi}{\mathrm{d}t} = 2\sqrt{\dfrac{\mu S}{R}}\sin\dfrac{\varphi}{2}$，故
\[
  \mathrm{d}t = \frac{1}{2}\sqrt{\frac{R}{\mu S}}\frac{\mathrm{d}\varphi}{\sin(\varphi/2)}.
\]
通过第二象限 $\varphi: \pi/2 \to \pi$：
\[
  t_2 - t_1 = \frac{1}{2}\sqrt{\frac{R}{\mu S}}\int_{\pi/2}^\pi \frac{\mathrm{d}\varphi}{\sin(\varphi/2)} = \sqrt{\frac{R}{\mu S}}\ln\tan\frac{\varphi}{4}\Big|_{\pi/2}^\pi \approx 0.88\sqrt{\frac{R}{\mu S}}.
\]

% === 例 11 ===
\item 重力与弹性力都是保守力。由机械能守恒（始末动能均为 0）：
\[
  mgy_{\max} = \frac{1}{2}ky_{\max}^2 \Rightarrow y_{\max} = \frac{2mg}{k}.
\]

% === 例 12 ===
\item 设给 $m_2$ 压力 $F$ 使弹簧压缩 $x_1 = F/k$；松手后机械能守恒，$m_1$ 离开桌面时弹簧拉伸 $x_2 = m_1 g/k$。$m_1$ 离开桌面的条件是弹簧恰拉伸 $x_2$。\\
取 $m_2$ 压后状态与 $m_1$ 离桌面状态间的机械能守恒（重力、弹力都是保守力）：
\[
  \frac{1}{2}k(x_0 + x_1)^2 = \frac{1}{2}kx_2^2 + m_2 g(x_0 + x_1 + x_2).
\]
代入 $x_0 = m_2 g/k$，$x_1 = F/k$，$x_2 = m_1 g/k$，解得
\[
  F = (m_1 + m_2)g.
\]

% === 例 13 ===
\item 选 $B$ 点为重力势能零点。$A$ 点（弹簧原长，水平位置）相对 $B$ 高 $R + R\cos 60° = 1.5R$。\\
初态（$A$）：$E_k = 0$，$E_{p弹} = 0$，$E_{p重} = mg\cdot 1.5R$。\\
末态（$B$）：$E_k$，$E_{p弹} = \dfrac{1}{2}kR^2$（$B$ 处弹簧伸长 $R$），$E_{p重} = 0$。\\
机械能守恒：$E_k + \dfrac{1}{2}kR^2 = \dfrac{3}{2}mgR$，故 $E_k = \dfrac{3}{2}mgR - \dfrac{1}{2}kR^2$。

% === 例 14 ===
\item 由机械能守恒（万有引力是保守力）：$\dfrac{1}{2}mv_0^2 - G\dfrac{M_e m}{R_e} = \dfrac{1}{2}mv^2 - G\dfrac{M_e m}{x}$。\\
代入 $v_0 = \sqrt{2GM_e/R_e}$，得 $v = \sqrt{2GM_e/x}$。\\
$\mathrm{d}t = \dfrac{\mathrm{d}x}{v} = \dfrac{1}{\sqrt{2GM_e}}\sqrt{x}\,\mathrm{d}x$。
\[
  t_1 = \int_{R_e}^{nR_e}\frac{\sqrt{x}\,\mathrm{d}x}{\sqrt{2GM_e}} = \frac{1}{\sqrt{2GM_e}}\cdot\frac{2}{3}x^{3/2}\Big|_{R_e}^{nR_e} = \frac{2R_e^{3/2}(n^{3/2}-1)}{3\sqrt{2GM_e}}.
\]

\end{enumerate}

\end{document}
